\section{图表与材料清单}

\subsection{已生成图表}

\begin{itemize}
  \item \code{5m\_vs\_6m} QMIX variants win-rate learning curves。
  \item \code{3s5z} QMIX variants win-rate learning curves。
  \item \code{5m\_vs\_6m} IQL/VDN/QMIX cross-algorithm AttnRes-L2 learning curves。
  \item \code{3s5z} IQL/VDN/QMIX cross-algorithm AttnRes-L2 learning curves（follow-up 完成后生成）。
  \item \code{5m\_vs\_6m} QMIX variants training wall-clock bar chart。
  \item AttnRes-L2 attention heatmaps（基于记录 attention weights 的 seed4/5 runs）。
  \item PyMARL/QMIX 训练框架图。
  \item 原始 RNN agent、AttnRes-RNN agent、depth-only control 结构图。
\end{itemize}

\subsection{当前实验完整性}

当前预设实验中，QMIX 主线和 \code{5m\_vs\_6m} 跨算法验证均已补齐 3 个 seeds。
VDN baseline 新增 run 为：

\begin{itemize}
  \item \code{5m\_vs\_6m vdn seed1}: Sacred run 100。
  \item \code{5m\_vs\_6m vdn seed2}: Sacred run 101。
  \item \code{5m\_vs\_6m vdn seed3}: Sacred run 102。
\end{itemize}

重新汇总后，\code{marl\_transfer\_missing\_or\_partial.csv} 中无预设缺失项。
本轮 follow-up 使用 GPU 0/2/3/4/5/6 补充 \code{3s5z} 跨算法验证，并将核心 QMIX、AttnRes-L2 和 depth-only control 扩展到 seed4/5；Full/Block 仍保持 3 seeds。

\subsection{后续可扩展参考方向}

\begin{itemize}
  \item MARL 与 value decomposition：可继续补充 IQL、VDN、QMIX 之后的代表性改进方法。
  \item SMAC benchmark：可补充 SMACv2 或更复杂地图上的近期评测工作。
  \item Attention in MARL：可继续扩展 communication、graph attention 和 transformer-based MARL。
  \item Residual 与 Transformer：可补充更多跨层聚合、PreNorm 和深层网络稳定性研究。
  \item RL stability：可补充 off-policy value learning 和部分可观测 recurrent agents 的稳定性分析。
\end{itemize}
